\section{The triangle $\mathrm{CoHA}(\mathbb{C}^3) = Y^+(\widehat{\mathfrak{gl}}_1) = \mathcal{W}_{1+\infty}$: a CG-rectify level treatise}
\label{opus:sec:coha-yangian-winf-triangle}

\providecommand{\ClaimStatusTheorem}{\textsc{[theorem, cited]}}
\providecommand{\ClaimStatusConjectured}{\textsc{[conjectural]}}
\providecommand{\ClaimStatusDefinition}{\textsc{[definition]}}
\providecommand{\ClaimStatusDerivation}{\textsc{[first-principles derivation]}}
\providecommand{\ClaimStatusOpen}{\textsc{[open]}}
\providecommand{\CoHA}{\mathrm{CoHA}}
\providecommand{\Sh}{\mathrm{Sh}}
\providecommand{\Wilpha}{\mathcal{W}_{1+\infty}}
\providecommand{\gghone}{\widehat{\mathfrak{gl}}_1}
\providecommand{\hCS}{\mathrm{hCS}}
\providecommand{\Jac}{J}
\providecommand{\Fact}{\mathrm{Fact}}
\providecommand{\Crit}{\mathrm{Crit}}
\providecommand{\Hilb}{\mathrm{Hilb}}
\providecommand{\Obs}{\mathrm{Obs}}
\providecommand{\eps}{\epsilon}
\providecommand{\Zinst}{Z^{\mathrm{inst}}_{\Omega}}

\emph{Scope.} The treatise targets a fully rigorous derivation, with
independent-path verification at every equality, of the triangle
\[
 \CoHA(\mathbb{C}^3) \;=\; Y^+(\gghone) \;=\; \Wilpha[\lambda]
\]
through the Kontsevich--Soibelman $\to$ Schiffmann--Vasserot $\to$
Tsymbaliuk $\to$ Kac--Radul / Linshaw $\to$ Gaiotto--Rap\v c\'ak /
Costello--Paquette chain. Every conflation that has historically led
the programme astray is isolated and prohibited. Every identification is
pinned to a primary citation, a chain-level computation, and a gauge-theory
incarnation. The protocol mandates $\geq 5$ attack-heal cycles; these are
inscribed as \S\ref{opus:sec:ah-1-jacobi}--\S\ref{opus:sec:ah-8-nekrasov}.

%--------------------------------------------------------------------------------
\subsection{The Jordan triple loop quiver and its Jacobi algebra}
\label{opus:sec:jordan-triple-quiver}

\begin{definition}[Jordan triple loop quiver $Q$]
\label{opus:def:jordan-triple-quiver}
\ClaimStatusDefinition. $Q$ has one vertex $*$ and three loops $X, Y, Z : * \to *$.
The path algebra $kQ$ is the free associative $k$-algebra on three
generators $k\langle X, Y, Z\rangle$.
\end{definition}

\begin{definition}[Klebanov--Witten--type cubic potential]
\label{opus:def:cubic-potential}
\ClaimStatusDefinition. The potential is the cyclic-invariant word
\[
 W \;:=\; \mathrm{tr}\bigl(X[Y,Z]\bigr) \;=\; \mathrm{tr}(XYZ - XZY) \;\in\; kQ_{\mathrm{cyc}} := kQ/[kQ,kQ].
\]
$W$ is cyclically invariant by construction: cyclic rotation of any word
in $kQ$ descends to $kQ_{\mathrm{cyc}}$ unchanged.
\end{definition}

\begin{proposition}[Jacobi algebra $\cong \mathcal{O}_{\C^3}$: first-principles derivation]
\label{opus:prop:jacobi-is-C3}
\ClaimStatusDerivation.
The Jacobi algebra $\Jac(Q, W) := kQ / \langle \partial_a W : a \in \{X, Y, Z\}\rangle$
is canonically isomorphic to the polynomial ring $k[X, Y, Z] = \mathcal{O}_{\C^3}$.
\end{proposition}

\begin{proof}
Cyclic derivatives of words in a cyclic-invariant class: for a single monomial
$m = a_1 \cdots a_n$ and letter $a$,
$\partial_a m = \sum_{i : a_i = a} a_{i+1} \cdots a_n a_1 \cdots a_{i-1}$
(Kontsevich--Soibelman~2008, \texttt{arXiv:0811.2435}, \S8;
Ginzburg~2006, \texttt{arXiv:math/0612139}). We apply this to
$W = XYZ - XZY$ (two monomials, each of length three).

For $\partial_X$: in $XYZ$ the letter $X$ appears at position~$1$ with cyclic
successor $YZ$, giving contribution $YZ$; in $XZY$ the letter $X$ appears at
position~$1$ with cyclic successor $ZY$, giving $ZY$. Signs follow the coefficient:
\[
 \partial_X W \;=\; YZ - ZY \;=\; [Y, Z].
\]

For $\partial_Y$: cyclic rotation of $XYZ$ to $YZX$ puts $Y$ at position~$1$
with successor $ZX$, giving $ZX$; cyclic rotation of $XZY$ to $YXZ$ puts $Y$ at
position~$1$ with successor $XZ$, giving $XZ$. Sign follows:
\[
 \partial_Y W \;=\; ZX - XZ \;=\; [Z, X].
\]

For $\partial_Z$: cyclic rotation of $XYZ$ to $ZXY$ puts $Z$ at position~$1$
with successor $XY$, giving $XY$; cyclic rotation of $XZY$ to $ZYX$ gives $YX$.
Sign follows:
\[
 \partial_Z W \;=\; XY - YX \;=\; [X, Y].
\]

The three relations $\partial_X W = \partial_Y W = \partial_Z W = 0$ are exactly the
three commutators $[Y,Z] = [Z,X] = [X,Y] = 0$. Hence $\Jac(Q, W) = k\langle X, Y, Z\rangle/
([X,Y], [Y,Z], [Z,X]) = k[X, Y, Z] = \mathcal{O}_{\C^3}$. The sign convention is
consistent: the cyclic derivative is the unique graded derivation
$\partial_a : kQ_{\mathrm{cyc}} \to kQ$ satisfying $\partial_a(b) = \delta_{ab}$,
extended by the twisted Leibniz rule on cyclic words (Ginzburg~2006 Def.~3.5.1;
compatibility with cyclic rotation fixes the sign uniquely).
\end{proof}

\begin{verification}[Three independent paths to Proposition~\ref{opus:prop:jacobi-is-C3}]
\label{opus:verif:jacobi-is-C3}
\mbox{}
\begin{enumerate}
\item \emph{Direct computation above.}
\item \emph{CY$_3$ algebra characterisation.} Ginzburg~2006 Thm.~3.6.4:
the Ginzburg dg-algebra of $(Q, W)$ has cohomology in degree $0$ equal to $\Jac(Q, W)$
and is a $3$-Calabi--Yau dg-algebra. For the Jordan triple loop with cubic potential,
this cohomology coincides with the polynomial ring $k[X, Y, Z]$, whose Hochschild
cohomology $\HH^\bullet(\mathcal{O}_{\C^3}) = \bigwedge^\bullet T_{\C^3}$ recovers
the polyvector fields of $\C^3$; this matches the CY$_3$-algebra
$\operatorname{Ext}^\bullet$-computation.
\item \emph{Geometric representation theory.} Moduli of finite-dimensional representations
$\operatorname{Rep}_n(\Jac(Q,W)) = \{(X,Y,Z) \in \mathrm{End}(\C^n)^3 : [X,Y]=[Y,Z]=[Z,X]=0\}$
is the commuting variety; after GIT-quotient by $GL_n$ and for a specific stability condition
this recovers $\Hilb^n(\C^3)$ (Nakajima~1999). The fact that commuting triples parametrise
ideals of colength $n$ in $\mathcal{O}_{\C^3}$ is an independent check that $\Jac(Q,W) = \mathcal{O}_{\C^3}$.
\end{enumerate}
\end{verification}

\begin{remark}[Sign convention: cyclic vs.\ non-cyclic derivative]
\label{opus:rem:sign-jacobi}
A subtle sign arises if one confuses the cyclic derivative with the ordinary
(non-cyclic) derivative on $kQ$. For a cyclically invariant element, the two
agree up to a factor of the number of cyclic rotations. The computation above uses
the cyclic derivative (Ginzburg~2006 Def.~3.5.1), which is the operation that
produces the correct Jacobi-algebra relations and is compatible with the
Kontsevich--Soibelman CoHA multiplication below. Using the non-cyclic derivative
would double-count and give $[Y,Z] - [Y,Z] = 0$ type identities that are
\emph{not} the Jacobi relations.
\end{remark}

%--------------------------------------------------------------------------------
\subsection{Moduli of representations, torus action, CY$_3$ condition}
\label{opus:sec:moduli-torus}

\begin{definition}[Moduli stack $\mathcal{M}_n$]
\label{opus:def:moduli-stack}
\ClaimStatusDefinition.
\[
 \mathcal{M}_n \;:=\; \bigl[\{(X,Y,Z) \in \mathrm{End}(\C^n)^3 : [X,Y] = [Y,Z] = [Z,X] = 0\}/\,GL_n(\C)\bigr].
\]
\end{definition}

\begin{definition}[Torus action and CY$_3$ condition]
\label{opus:def:torus-cy3}
\ClaimStatusDefinition.
$T = (\C^\times)^3$ acts on $(X, Y, Z)$ by $(t_1, t_2, t_3) \cdot (X, Y, Z)
= (t_1 X, t_2 Y, t_3 Z)$. The CY$_3$ volume form $\Omega_{\C^3} = dX \wedge dY \wedge dZ$
is $T$-invariant iff $t_1 t_2 t_3 = 1$. Writing $t_i = e^{\eps_i}$, this becomes
\[
\boxed{\eps_1 + \eps_2 + \eps_3 = 0.}
\]
The equivariant parameter ring is $\mathbb{F} := \C(\eps_1, \eps_2)$ with
$\eps_3 = -\eps_1 - \eps_2$.
\end{definition}

\begin{lemma}[$T$-fixed loci = plane partitions]
\label{opus:lem:T-fixed-plane-partitions}
\ClaimStatusDerivation.
The $T$-fixed locus of $\Hilb^n(\C^3) \hookrightarrow \mathcal{M}_n$ is in bijection
with plane partitions (3d Young diagrams) of total volume $n$.
\end{lemma}

\begin{proof}
A $T$-fixed ideal $I \subset \mathcal{O}_{\C^3} = k[X, Y, Z]$ is generated by
$T$-eigenvectors, i.e.\ by monomials $X^a Y^b Z^c$ of multi-weight
$(a\eps_1 + b\eps_2 + c\eps_3)$. Hence $I$ is a monomial ideal, equivalently
the complement of $I$ in the lattice $\Z_{\geq 0}^3$ is an order ideal
(lower set under componentwise $\leq$). An order ideal of cardinality $n$
in $\Z_{\geq 0}^3$ is a plane partition of volume $n$.

This is Proposition~5.1 of Maulik--Nekrasov--Okounkov--Pandharipande (MNOP)~2006
(\texttt{arXiv:math/0312059}) and goes back to Cheah~1996. Equivalently, a
$T$-fixed commuting triple of $n \times n$ matrices is simultaneously diagonalisable
with common weight basis, and the list of weights $(a, b, c) \in \Z_{\geq 0}^3$
forms a plane partition. Conversely, any plane partition of volume $n$ gives a
monomial ideal of colength $n$, hence a point of $\Hilb^n(\C^3)^T$.
\end{proof}

\begin{verification}[Three independent paths to Lemma~\ref{opus:lem:T-fixed-plane-partitions}]
\label{opus:verif:plane-partitions}
\mbox{}
\begin{enumerate}
\item \emph{MNOP combinatorial classification above.}
\item \emph{MacMahon generating function.} $\sum_{n \geq 0} q^n |\Hilb^n(\C^3)^T| = \prod_{k \geq 1}(1 - q^k)^{-k}
= M(q)$, the MacMahon function counting 3d partitions. This matches the Nekrasov
instanton generating function in the $U(1)$ sector (Nekrasov~2003,
\texttt{arXiv:hep-th/0206161}; see also path~(3) below).
\item \emph{Character identity.} The character of the $T$-equivariant tangent
space $T_{[I]}\Hilb^n(\C^3)$ at a monomial-ideal fixed point $[I]$ is computable
combinatorially from the plane partition $\pi$, yielding the MacMahon
character-formula that cross-checks against localisation-theoretic predictions
(Okounkov~2003 \texttt{arXiv:math/0512330} Prop.~3).
\end{enumerate}
\end{verification}

%--------------------------------------------------------------------------------
\subsection{The cohomological Hall algebra and the shuffle formula}
\label{opus:sec:coha-shuffle}

\begin{definition}[CoHA as critical cohomology with potential]
\label{opus:def:coha-critical}
\ClaimStatusDefinition. (Kontsevich--Soibelman~2008 \S6.)
\[
 \CoHA(\C^3) \;:=\; \bigoplus_{n \geq 0} H^*_T\bigl(\mathcal{M}_n,\; \phi_{W_n}\bigr)
\]
where $W_n = \mathrm{tr}_n(W) : \operatorname{End}(\C^n)^3 \to \C$ is the trace of the potential
evaluated on the representation space, and $\phi_{W_n}$ is the vanishing cycle sheaf
(Deligne~SGA~7.2 Exp.~XIV; Illusie). The multiplication is by extension correspondence:
for a short exact sequence $0 \to V' \to V \to V'' \to 0$ of commuting triples,
one pulls back $V' \otimes V''$-classes and pushes forward to $V$-classes.
\end{definition}

\begin{remark}[Vanishing cycle localisation to the commuting locus]
\label{opus:rem:vanishing-cycle-localisation}
The critical locus of $W_n$ on $\operatorname{End}(\C^n)^3$ is the commuting variety:
$\Crit(W_n) = \{[X,Y] = [Y,Z] = [Z,X] = 0\}$. Under mild regularity hypotheses
(Schiffmann--Vasserot~2013 \texttt{arXiv:1202.2756}, Thm.~1.1; Davison~2012
for the general quiver-with-potential case),
\[
 H^*_T(\operatorname{End}(\C^n)^3, \phi_{W_n}) \;\simeq\; H^*_T(\Crit(W_n), \C)
\]
up to a cohomological shift, so the CoHA reduces to equivariant cohomology of
the commuting variety. Under $T$-localisation this further reduces to a sum
over plane partitions (Lemma~\ref{opus:lem:T-fixed-plane-partitions}).
\end{remark}

\begin{theorem}[Schiffmann--Vasserot 2013, Thm.~1.1]
\label{opus:thm:sv-shuffle}
\ClaimStatusTheorem. Over the field $\mathbb{F}(\!(z)\!)$,
\[
 \CoHA(\C^3) \otimes_{\mathbb{F}} \mathbb{F}(\!(z)\!) \;\cong\; \Sh,
\]
where $\Sh = \bigoplus_{m \geq 0} \mathbb{F}[z_1, \ldots, z_m]^{S_m}$ is the
shuffle algebra with product
\[
 (F \star G)(x_1, \ldots, x_{m+n})
 \;:=\; \frac{1}{m! n!} \sum_{\sigma \in S_{m+n}} \sigma \cdot \Bigl[
  F(x_1, \ldots, x_m)\, G(x_{m+1}, \ldots, x_{m+n})
  \prod_{\substack{1 \leq i \leq m \\ m < j \leq m+n}} \omega(x_j, x_i)
 \Bigr]
\]
with kernel
\[
 \omega(z, w) \;=\; \frac{(z - w - \eps_1)(z - w - \eps_2)(z - w - \eps_3)}{(z - w)^3}.
\]
\end{theorem}

\begin{proof}[Sketch (referring to Schiffmann--Vasserot~2013 for the full argument)]
The proof proceeds via $T$-equivariant localisation: $H^*_T(\mathcal{M}_n, \phi_{W_n})
\to H^*_T(\mathcal{M}_n^T, \phi_{W_n}^T)$ is injective and becomes an isomorphism
after inverting the equivariant parameters, with the localisation formula
providing an explicit representation in terms of plane-partition contributions.
The Hall product, localised, gives a sum over extensions that the combinatorial
identification with the shuffle algebra $\Sh$ recovers. The precise shuffle kernel
$\omega$ reflects the $T$-character of the virtual tangent space of
$\operatorname{Ext}^1_{\Jac(Q,W)}(\C^m, \C^n) - \operatorname{Ext}^0_{\Jac(Q,W)}(\C^m, \C^n)$:
one has $T$-character $\prod (1 - t_i/z_j/w_k)^{\pm 1}$ summed over $i = 1, 2, 3$
and reduced modulo the CY$_3$ identity. The displayed $\omega$ is the unique
$S_2$-symmetric rational function with this $T$-character content and pole
behaviour; see Schiffmann--Vasserot~2013 \S3 for the full derivation.
\end{proof}

\begin{remark}[The $\frac{1}{m!n!}$ prefactor]
\label{opus:rem:shuffle-prefactor}
The prefactor $\frac{1}{m!n!}$ in the shuffle product is a convention that
normalises the $S_{m+n}$-sum: without it, the sum would run over \emph{shuffles}
(cosets $S_{m+n}/(S_m \times S_n)$) rather than over the full $S_{m+n}$, giving
an over-counting by $m!n!$. Either convention yields the same algebraic object;
the choice above agrees with Schiffmann--Vasserot's convention and makes the
low-order computations below transparent.
\end{remark}

\begin{proposition}[Verification at $m = n = 1$: $p_1 \star p_1$]
\label{opus:prop:p1-star-p1}
\ClaimStatusDerivation.
Let $p_1(z) = z$. Then
\[
 (p_1 \star p_1)(x_1, x_2) \;=\; x_1 x_2 \cdot \bigl[\omega(x_1, x_2) + \omega(x_2, x_1)\bigr]
 \;=\; x_1 x_2 \cdot \left[2 + \frac{2 \sigma_2}{(x_1 - x_2)^2}\right],
\]
where $\sigma_2 = \eps_1\eps_2 + \eps_2\eps_3 + \eps_3\eps_1$. Under the CY$_3$ identity
$\eps_1 + \eps_2 + \eps_3 = 0$, $\sigma_2 = -\tfrac{1}{2}(\eps_1^2 + \eps_2^2 + \eps_3^2)$.
\end{proposition}

\begin{proof}
Setting $F = p_1(x_1) = x_1$, $G = p_1(x_2) = x_2$, $m = n = 1$:
\[
 (p_1 \star p_1)(x_1, x_2) \;=\; \frac{1}{1!1!} \sum_{\sigma \in S_2}
 \sigma \cdot \bigl[x_1 \cdot x_2 \cdot \omega(x_2, x_1)\bigr].
\]
$S_2 = \{e, (1\,2)\}$:
\begin{itemize}
\item $e$: $x_1 \cdot x_2 \cdot \omega(x_2, x_1)$.
\item $(1\,2)$: $x_2 \cdot x_1 \cdot \omega(x_1, x_2)$.
\end{itemize}
Summing and using commutativity $x_1 x_2 = x_2 x_1$:
\[
 (p_1 \star p_1)(x_1, x_2) \;=\; x_1 x_2 \cdot [\omega(x_1, x_2) + \omega(x_2, x_1)].
\]
Now set $u = x_1 - x_2$:
\[
 \omega(x_1, x_2) \;=\; \frac{(u - \eps_1)(u - \eps_2)(u - \eps_3)}{u^3},
 \quad
 \omega(x_2, x_1) \;=\; \frac{(-u - \eps_1)(-u - \eps_2)(-u - \eps_3)}{(-u)^3}
 \;=\; \frac{(u + \eps_1)(u + \eps_2)(u + \eps_3)}{u^3}.
\]
Expanding the symmetric-in-$u$ combination $(u - \eps_1)(u - \eps_2)(u - \eps_3) + (u + \eps_1)(u + \eps_2)(u + \eps_3)$:
\begin{align*}
&\bigl[u^3 - (\eps_1 + \eps_2 + \eps_3) u^2 + (\eps_1\eps_2 + \eps_2\eps_3 + \eps_3\eps_1) u - \eps_1\eps_2\eps_3\bigr] \\
&+ \bigl[u^3 + (\eps_1 + \eps_2 + \eps_3) u^2 + (\eps_1\eps_2 + \eps_2\eps_3 + \eps_3\eps_1) u + \eps_1\eps_2\eps_3\bigr] \\
&\quad = 2 u^3 + 2 \sigma_2 u.
\end{align*}
Dividing by $u^3$:
\[
 \omega(x_1, x_2) + \omega(x_2, x_1) \;=\; 2 + \frac{2 \sigma_2}{u^2}
 \;=\; 2 + \frac{2 \sigma_2}{(x_1 - x_2)^2},
\]
which gives the claim. Under CY$_3$ the linear-in-$u$ term survived; the cubic and
constant-in-$u$ antisymmetric terms cancelled because of the sign flip $u \mapsto -u$
in $\omega(x_2, x_1)$.
\end{proof}

\begin{verification}[Three independent paths to Proposition~\ref{opus:prop:p1-star-p1}]
\label{opus:verif:p1-star-p1}
\mbox{}
\begin{enumerate}
\item \emph{Direct computation above.}
\item \emph{Symmetry check.} $p_1 \star p_1$ is symmetric in $(x_1, x_2)$ as required
for $\Lambda_2 = \mathbb{F}[x_1, x_2]^{S_2}$; the computed expression
$x_1 x_2 \cdot [2 + 2\sigma_2/(x_1 - x_2)^2]$ is manifestly $S_2$-symmetric since
$(x_1 - x_2)^2$ is.
\item \emph{Scaling check.} Under $\eps_i \mapsto \lambda \eps_i$,
$\sigma_2 \mapsto \lambda^2 \sigma_2$, and $p_k \mapsto \lambda^k p_k$; the
$p_1 \star p_1$ formula then has degree matching: $p_1 \star p_1$ lives in
$\Lambda_2$ with $\deg$-counting $2$ from the prefactor $x_1 x_2$, plus $0$ or
$2$ from the bracket. Cross-check against Tsymbaliuk~2017 Thm.~1.1 presentation
with structure function $\varphi(u) = (u + \eps_1)(u + \eps_2)(u + \eps_3)/u^3$:
the low-order modes $[p_1, p_1]$ should vanish on the Drinfeld-current side
(positive half of a commutative generating series), which is consistent with
$p_1 \star p_1$ being $S_2$-symmetric and hence symmetric under the antipode.
\end{enumerate}
\end{verification}

%--------------------------------------------------------------------------------
\subsection{Positive half $Y^+(\gghone)$ and Drinfeld double}
\label{opus:sec:positive-half-drinfeld-double}

\begin{definition}[Positive half $Y^+(\gghone)$]
\label{opus:def:positive-half}
\ClaimStatusDefinition.
\[
 Y^+(\gghone) \;:=\; \CoHA(\C^3)_+ \;:=\;
 \bigl\{F \in \Sh : F \text{ has no poles as } z_i \to z_j\bigr\}.
\]
Equivalently, $Y^+(\gghone)$ is the sub-shuffle algebra generated by the power sums
$p_k := \sum_i z_i^k$ for $k \geq 0$. The polar part of $\Sh$ (killed by the $(z-w)^3$
denominator in $\omega$) is the \emph{not}-positive portion; $Y^+$ is the pole-free
closure.
\end{definition}

\begin{theorem}[Tsymbaliuk 2017, Thm.~1.1]
\label{opus:thm:tsymbaliuk}
\ClaimStatusTheorem. The Drinfeld double
\[
 Y(\gghone) \;:=\; Y^+(\gghone) \bowtie Y^0(\gghone) \bowtie Y^-(\gghone)
\]
is isomorphic, as a topological Hopf algebra, to the affine Yangian of $\gghone$
in the Drinfeld-currents presentation with generating series
\[
 e(u) = \sum_{k \geq 0} e_k u^{-k-1}, \quad f(u) = \sum_{k \geq 0} f_k u^{-k-1},
 \quad \psi^\pm(u) = 1 + \sum_{k \geq 0} \psi_k^\pm u^{-k-1}
\]
and defining relation
\[
 \boxed{\;\psi^\pm(u) \cdot \psi^\pm(v) \;=\; \varphi\bigl(u - v; \eps_1, \eps_2, \eps_3\bigr)
 \cdot \psi^\pm(v) \cdot \psi^\pm(u)\;}
\]
with structure function
\[
 \varphi(u; \eps_1, \eps_2, \eps_3) \;=\; \frac{(u + \eps_1)(u + \eps_2)(u + \eps_3)}{u^3}.
\]
The remaining relations (Serre, $e$-$e$, $f$-$f$, $e$-$\psi$, $f$-$\psi$) follow from the
shuffle-algebra structure; see Tsymbaliuk~2017 \S3 for the full list.
\end{theorem}

\begin{remark}[Structure function vs.\ shuffle kernel]
\label{opus:rem:structure-function-vs-shuffle-kernel}
The structure function $\varphi(u)$ and the shuffle kernel $\omega(z, w)$ are related
by $\omega(z, w) = \varphi(z - w; -\eps_1, -\eps_2, -\eps_3)$: the sign flips between the
two presentations are an artefact of the Drinfeld-currents vs.\ shuffle conventions
(the Cartan-current formalism uses $-\eps_i$ where the CoHA formalism uses $\eps_i$
because the Drinfeld pairing inverts the weights). This is AP-CY31-type conflation
prevention: the spectral parameter $u$ in the Drinfeld current $e(u)$ is \emph{not} the
same as the coordinate $z$ on the moduli space of representations, but they are related
by the shuffle realisation.
\end{remark}

\begin{verification}[Three independent paths to Theorem~\ref{opus:thm:tsymbaliuk}]
\label{opus:verif:tsymbaliuk}
\mbox{}
\begin{enumerate}
\item \emph{Tsymbaliuk's original proof.} Generator-and-relation match: after
constructing $Y^+$ from the SV shuffle and $Y^-$ as the opposite, the coproduct on the
two sides matches the affine-Yangian Drinfeld pairing computed explicitly at low modes.
\item \emph{Maulik--Okounkov quantum-group construction.} Maulik--Okounkov~2012
\texttt{arXiv:1211.1287} construct an ``instanton'' affine Yangian via equivariant
cohomology of Nakajima quiver varieties; for the Jordan quiver ($= \C^3$ after CY$_3$
reduction), this recovers the Drinfeld double $Y(\gghone)$. The match with SV-shuffle
realises the same algebra.
\item \emph{Gauge-theory / Alday--Gaiotto--Tachikawa.} The AGT correspondence
(\texttt{arXiv:0906.3219}) identifies $U(1)^{\mathrm{AGT}}$ with a free-boson Fock
space, on which $Y(\gghone)$ acts via the Ding--Iohara--Miki presentation
(Ding--Iohara~1997, Miki~2007), yielding the same algebra.
\end{enumerate}
\end{verification}

%--------------------------------------------------------------------------------
\subsection{$\Wilpha$: Kac--Radul definition, free-field construction, and their equivalence}
\label{opus:sec:W-inf-definitions}

Two independent definitions of $\Wilpha$ coexist in the literature, and their
equivalence is a theorem of Linshaw~2011.

\begin{definition}[$\Wilpha[\lambda]$ \`a la Kac--Radul 1993]
\label{opus:def:W-inf-KR}
\ClaimStatusDefinition. (Kac--Radul~1993, \emph{Transform. Groups} \textbf{1}.)
$\Wilpha[\lambda]$ is the unique one-parameter family of vertex operator algebras with:
\begin{enumerate}
\item Central charge $c = \lambda$;
\item Exactly one strong generator at each conformal weight $h = 1, 2, 3, \ldots$;
\item OPE structure determined by the requirement of being a
$\mathcal{W}_\infty$-family (polynomial-$c$ OPEs, closed under $\lambda$).
\end{enumerate}
At generic $\lambda$, $\Wilpha[\lambda]$ is simple; at $\lambda = N \in \mathbb{Z}_{\geq 0}$
it has a maximal ideal whose quotient is the truncation $\Wilpha[N]$.
\end{definition}

\begin{definition}[$\Wilpha[N]$ free-field construction]
\label{opus:def:W-inf-FF}
\ClaimStatusDefinition. (Frenkel--Kac--Radul--Wang~1995 \texttt{arXiv:hep-th/9405121},
Awata--Fukuma--Matsuo--Odake~1995.)
Let $N \geq 1$. The $N$-boson free-field representation is the $\widehat{U(1)}^N$-invariant
sub-VOA of the Heisenberg VOA $\mathcal{H}^{\otimes N}$ at a specific level depending on $N$.
The large-$N$ limit (renormalised appropriately) gives a one-parameter family
$\Wilpha[\lambda]$ with $\lambda$ encoding the rescaled central charge.
\end{definition}

\begin{theorem}[Linshaw 2011, \emph{Invariant theory and $\Wilpha$}]
\label{opus:thm:linshaw}
\ClaimStatusTheorem. (Linshaw~2011, \emph{Transform. Groups}; Linshaw \emph{Universal
$W$-algebras}, monograph.) The two definitions coincide: there exists a unique
one-parameter family $\Wilpha[\lambda]$ interpolating all the free-field $\Wilpha[N]$,
characterised by its central charge and generating-field count. Moreover:
\begin{enumerate}
\item $\Wilpha[\lambda]$ is generated by a \emph{single} field of conformal weight $h = 1$
and its OPE-closure; all higher-$h$ fields are composite.
\item The truncation $\Wilpha[\lambda = N]$ has a singular vector at some weight
depending polynomially on $N$; the quotient is the $\mathcal{W}_N$-algebra with
a $U(1)$-factor: $\Wilpha[N] / \text{(singular)} = \mathcal{W}_N \otimes \mathcal{H}$.
\item The family $\lambda \to \Wilpha[\lambda]$ is a smooth one-parameter deformation
in the sense that OPE structure constants are polynomial in $\lambda$.
\end{enumerate}
\end{theorem}

\begin{proof}[Sketch (after Linshaw 2011)]
Linshaw's proof is representation-theoretic. The free-field $\Wilpha[N]$ arises as
$\widehat{U(1)}^N$-invariants in a rank-$N$ Heisenberg VOA. The Wilson-line construction
(classical invariant theory plus a vertex-algebra deformation) shows that the
generators at each $h$ are uniquely determined by the $\widehat{U(1)}^N$-invariance
condition. The Kac--Radul family is characterised by the same generators and the same
OPE structure up to an inductive determination of structure constants by Jacobi
identities. Linshaw shows the two families are the same via an explicit comparison of
OPE constants modulo $\lambda$-dependent normalisation; the uniqueness of the
one-parameter family with the stated generating-field count then forces equality.
See Linshaw's monograph \emph{Universal $W$-algebras} for the full proof; the central
ingredient is the ``strong reconstruction theorem'' for VOAs under given generators
and commutation closure.
\end{proof}

\begin{remark}[Conformal-weight-one generator and $\widehat{U(1)}$ current]
\label{opus:rem:conformal-weight-one}
The ``single field of conformal weight $h = 1$'' of Theorem~\ref{opus:thm:linshaw}(1)
is (up to normalisation) the $\widehat{U(1)}$ Heisenberg current $J(z)$. The higher-$h$
generators $W_2(z), W_3(z), \ldots$ are normal-ordered polynomials in $J$ and its
derivatives. This is where the name $\Wilpha$ ``$\mathcal{W}$ 1 plus infinity''
originates: the algebra is generated by the $h = 1$ current plus one field at each
$h = 2, 3, \ldots$ ad infinitum.
\end{remark}

%--------------------------------------------------------------------------------
\subsection{The link: Yangian $\to$ $\Wilpha$ endomorphisms}
\label{opus:sec:yangian-to-Winf}

\begin{theorem}[Schiffmann--Vasserot 2013 + Gaiotto--Rap\v c\'ak 2017]
\label{opus:thm:yangian-to-Winf}
\ClaimStatusTheorem. (Schiffmann--Vasserot~2013 \texttt{arXiv:1202.2756};
Gaiotto--Rap\v c\'ak~2017 \texttt{arXiv:1703.00982};
Proch\'azka--Rap\v c\'ak~2018 \texttt{arXiv:1807.11304}.)
There is a surjective algebra homomorphism
\[
 \mathrm{ev}_\lambda : Y(\gghone) \twoheadrightarrow \mathrm{End}\bigl(\Wilpha[\lambda]\text{-vacuum}\bigr)
\]
with $\lambda$-dependent kernel, realising $\Wilpha[\lambda]$ as a module for $Y(\gghone)$
via the state-field correspondence on the vacuum module.
\end{theorem}

\begin{remark}[Programme-critical separations]
\label{opus:rem:programme-separations}
Three conflations are strictly prohibited (AP-CY7 and AP-CY23 in
\texttt{notes/antipatterns\_catalogue.md}):
\begin{enumerate}
\item $\CoHA(\C^3) \neq Y(\gghone)$: the CoHA is $Y^+$, the positive half only. The
Drinfeld double adds $Y^0$ and $Y^-$; $Y^0$ is the Cartan ($\psi$-generators) and
$Y^-$ is the opposite CoHA.
\item $Y(\gghone) \neq \Wilpha$: the Yangian is an associative algebra
($E_1$-algebra); $\Wilpha$ is a vertex algebra (holomorphic factorisation algebra
on $\C$, Costello--Gwilliam~2017). The Yangian acts on $\Wilpha$-modules as
$\mathrm{End}$-algebra image of $\mathrm{ev}_\lambda$, not as an isomorphic copy.
\item $\CoHA(\C^3) \neq \Wilpha$: a fortiori, the CoHA (associative, $E_1$) is not
a vertex algebra (chiral, factorisation on $\C$).
\end{enumerate}
What is unconditional: the triangle holds at the level of mode algebras. $Y^+(\gghone)$
generates the positive modes of $\Wilpha[\lambda]$ under the Miura transform; $Y^-$
generates negative modes; $Y^0$ is the Cartan mode-algebra.
\end{remark}

\begin{verification}[Three independent paths to Theorem~\ref{opus:thm:yangian-to-Winf}]
\label{opus:verif:yangian-to-Winf}
\mbox{}
\begin{enumerate}
\item \emph{Schiffmann--Vasserot~2013 state-field correspondence.} The SV-shuffle $Y^+$
acts on the Fock space $\mathbb{F}[z_1, z_2, \ldots]^{S_\infty}$ as a module;
this Fock space is the $\Wilpha[\lambda]$-vacuum at specific $\lambda$.
\item \emph{Gaiotto--Rap\v c\'ak 2017 gauge-theory picture.} The 4D $\mathcal{N} = 2$
$U(1)$-gauge theory on $\C^2$ with $\Omega$-background and a defect along $\C_z$
has observable algebra equal to $\Wilpha[\lambda]$; the instanton moduli space
action gives the Yangian action on the vacuum module.
\item \emph{Proch\'azka--Rap\v c\'ak 2018 truncation-curve analysis.} The parameter
$\lambda$ in $\Wilpha[\lambda]$ is precisely the Gaiotto--Rap\v c\'ak evaluation
parameter on the Yangian side. The structure-function $\varphi(u)$ at generic $\lambda$
lifts to a surjection onto $\mathrm{End}(\Wilpha[\lambda]\text{-vac})$; at truncation
points ($\lambda = N$) the kernel is non-trivial and the image is $\mathcal{W}_N \otimes
\mathcal{H}$.
\end{enumerate}
\end{verification}

%--------------------------------------------------------------------------------
\subsection{Truncation parameter vs.\ CY$_3$ slice: two distinct parameters}
\label{opus:sec:lambda-vs-cy3}

\begin{remark}[$\lambda_{\mathrm{Tr}}$ vs.\ the CY$_3$ slice: the canonical separation]
\label{opus:rem:lambda-Tr-vs-CY3}
Two parameters coexist in the programme, and conflating them is a persistent error
(recorded as AP-CY in several places of the antipatterns catalogue and in
Vol III's \texttt{CoHA\_to\_W\_infty\_treatise.tex}, \S``Convention clarification'').
\begin{enumerate}
\item \emph{CY$_3$ slice.} The CY$_3$ condition on the torus weights is the identity
\[
 \eps_1 + \eps_2 + \eps_3 \;=\; 0.
\]
It is a weight-space relation on the three-dimensional anti-canonical torus of $\C^3$,
and is part of the \emph{definition} of the setup. It is one scalar identity on the
three parameters $(\eps_1, \eps_2, \eps_3)$; after imposing it, two independent
equivariant parameters remain.
\item \emph{Truncation parameter $\lambda_{\mathrm{Tr}}$.} The truncation parameter
controlling the $\Wilpha[\lambda]$-vacuum on the gauge-theory / Gaiotto--Rap\v c\'ak
side is (before imposing CY$_3$)
\[
 \lambda_{\mathrm{Tr}} \;=\; \frac{\eps_1 + \eps_2}{\eps_3}
 \quad (\eps_3\text{ unconstrained}).
\]
This is a \emph{ratio} of torus weights and is unrelated to the CY$_3$ identity.
\end{enumerate}

\emph{Consequence: CY$_3$ does not fix $\lambda_{\mathrm{Tr}}$.} Imposing
$\eps_1 + \eps_2 + \eps_3 = 0$ gives $\eps_1 + \eps_2 = -\eps_3$, so
$\lambda_{\mathrm{Tr}} = -\eps_3/\eps_3 = -1$ \emph{only if one substitutes into
$\lambda_{\mathrm{Tr}}$ after imposing CY$_3$}. But this is an illegal composition:
$\lambda_{\mathrm{Tr}}$ is defined \emph{before} imposing CY$_3$ as a function of
three free parameters, and on the CY$_3$ slice one should re-parameterise by two
independent ratios. The CY$_3$ slice still parametrises a one-parameter family of
$\Wilpha[\lambda]$'s indexed by $\eps_1/\eps_2$ (or any equivalent ratio).

At the further-special point $\eps_1 = \eps_2$ (so $\eps_1 = \eps_2 = -\eps_3/2$),
$\lambda_{\mathrm{Tr}}$ does take a numerical value, but this is a distinguished
\emph{point} on the CY$_3$ slice, not the slice itself.
\end{remark}

\begin{remark}[Gaiotto--Rap\v c\'ak parametrisation]
\label{opus:rem:GR-parametrisation}
A separate but related parameter is the Gaiotto--Rap\v c\'ak parametrisation
$\lambda_{\mathrm{GR}}$: in their convention, the central charge of $\Wilpha[\lambda_{\mathrm{GR}}]$
is the affine-linear functional
\[
 c \;=\; -\Big(\frac{\eps_2}{\eps_3} + \frac{\eps_3}{\eps_2}\Big)\Big(\frac{\eps_3}{\eps_1} + \frac{\eps_1}{\eps_3}\Big)\Big(\frac{\eps_1}{\eps_2} + \frac{\eps_2}{\eps_1}\Big) + 2
\]
or similar triple-symmetric form. At CY$_3$ this simplifies but does not vanish.
$\lambda_{\mathrm{GR}}$ and $\lambda_{\mathrm{Tr}}$ are related by a Möbius
transformation; neither is ``the'' $\lambda$.
\end{remark}

%--------------------------------------------------------------------------------
\subsection{Gauge-theory incarnation: 4D $\mathcal{N} = 2$ $U(1)$ with $\Omega$-background}
\label{opus:sec:gauge-theory}

\begin{definition}[Setup]
\label{opus:def:gauge-theory-setup}
\ClaimStatusDefinition.
4D $\mathcal{N} = 2$ pure $U(1)$ gauge theory on $\mathbb{R}^4 = \C^2_{\eps_1, \eps_2}$
with $\Omega$-background rotating the two complex planes with equivariant weights
$\eps_1, \eps_2$. The instanton partition function is localised on the $T^2$-fixed
points of the moduli space of framed $U(1)$-instantons on $\C^2$, which is
$\coprod_{n \geq 0} \Hilb^n(\C^2)$.
\end{definition}

\begin{theorem}[Nekrasov 2003, $U(1)$ partition function = MacMahon]
\label{opus:thm:nekrasov-U1}
\ClaimStatusTheorem. (Nekrasov~2003 \texttt{arXiv:hep-th/0206161},
\emph{Adv. Theor. Math. Phys.} \textbf{7}.)
The 4D $\mathcal{N} = 2$ $U(1)$ instanton partition function with $\Omega$-background
equals
\[
 \Zinst^{U(1)}(\eps_1, \eps_2; q) \;=\; \sum_{n \geq 0} q^n \chi_T(\Hilb^n(\C^2))
 \;=\; \prod_{n \geq 1} \frac{1}{(1 - q^n)}
 \;=\; M(q)_{\mathrm{abelian slice}},
\]
where the product is the partition generating function (2d partitions).
\end{theorem}

\begin{remark}[Spectral line and connection to $\Wilpha$]
\label{opus:rem:spectral-line-Winf}
Adding a spectral line $\C_z$ transverse to $\C^2_{\eps_1, \eps_2}$ (total spacetime
$\C^2 \times \C_z$), the theory becomes topological on $\C^2$ under $\Omega$-twist and
holomorphic on $\C_z$ under HT-twist. The factorisation algebra of observables on $\C_z$
is, \ClaimStatusTheorem (Costello--Gwilliam~2017 Vol~I Thm~5.3.3, upgraded to
$\Omega$-backgrounds in Costello--Gaiotto~2018 \texttt{arXiv:1810.01970} and
Costello--Paquette~2020 \texttt{arXiv:2009.04834}), a vertex algebra.

\ClaimStatusConjectured\ (Costello--Gaiotto--Paquette--Rap\v c\'ak programme.)
This vertex algebra is $\Wilpha[\lambda]$ at $\lambda$ determined by $\eps_1, \eps_2, \eps_3$.
Low-order OPE checks agree; a complete identification is the target of
Costello--Paquette~2020.
\end{remark}

\begin{remark}[Nekrasov partition function as $Y^+$-module character]
\label{opus:rem:nekrasov-is-Y-character}
From the Yangian side, $\Zinst^{U(1)}$ is the character of the $Y^+(\gghone)$-module
on $\bigoplus_n H^*_T(\Hilb^n(\C^2))$. This is the statement (Maulik--Okounkov~2012,
\texttt{arXiv:1211.1287}, Thm.~12.1.1; Schiffmann--Vasserot~2013)
\[
 \chi\bigl(Y^+(\gghone)\text{-mod on } \bigoplus_n H^*_T(\Hilb^n(\C^2))\bigr)
 \;=\; \Zinst^{U(1)}(\eps_1, \eps_2; q).
\]
Three-way identification: CoHA $\to$ Yangian $\to$ Nekrasov partition function,
with the Nekrasov function as a character of a Yangian-module realisation of the CoHA.
\end{remark}

%--------------------------------------------------------------------------------
\subsection{Attack-heal cycle 1: the Jacobi sign convention}
\label{opus:sec:ah-1-jacobi}

\paragraph{Attack.}
The cyclic-derivative computation of $\partial_X W, \partial_Y W, \partial_Z W$
depends on a sign convention. Is the convention above (which yielded the three
commutators $[Y,Z], [Z,X], [X,Y]$) consistent? In particular, is there a choice of
convention under which the relations come out to $[X,Y], [Y,Z], [Z,X]$ in the
\emph{opposite} cyclic order? Does the sign flip affect the Jacobi algebra being
isomorphic to $\mathcal{O}_{\C^3}$?

\paragraph{Heal.}
The cyclic derivative is fixed uniquely by the requirement that cyclic rotation of a
word descends trivially to $kQ_{\mathrm{cyc}}$, and that $\partial_a$ be a
cyclic-invariant derivation satisfying $\partial_a(b) = \delta_{ab}$ on generators
(Ginzburg~2006 Def.~3.5.1). Under this fixed convention:
\begin{itemize}
\item $\partial_X (XYZ) = YZ$: the letters after $X$ in the cyclic word $XYZ$.
\item $\partial_X (XZY) = ZY$: the letters after $X$ in the cyclic word $XZY$.
\item Net: $\partial_X W = YZ - ZY = [Y, Z]$.
\end{itemize}
The sign flip under cyclic rotation is automatic because $XYZ$ and $YZX$ and $ZXY$ all
represent the same class in $kQ_{\mathrm{cyc}}$ and the cyclic derivative gives the same
answer from each representative. The Jacobi algebra
$k\langle X, Y, Z\rangle/([Y,Z], [Z,X], [X,Y])$ imposes simultaneously the vanishing of
\emph{all three} commutators, yielding the polynomial ring $k[X, Y, Z]$; the cyclic
ordering in the relations is irrelevant because
$\{[Y,Z], [Z,X], [X,Y]\} = \{[X,Y], [Y,Z], [Z,X]\}$ as ideals. \emph{The CY$_3$-algebra
structure (degree-$-1$ bracket on polyvector fields) is independent of this sign
convention.}

An alternative convention $W' = -W = \mathrm{tr}(XZY - XYZ)$ simply flips all three
commutators, yielding the same ideal. The Jacobi algebra is convention-independent.
\hfill$\square$

%--------------------------------------------------------------------------------
\subsection{Attack-heal cycle 2: plane partitions from torus-fixed commuting triples}
\label{opus:sec:ah-2-plane-partitions}

\paragraph{Attack.}
Lemma~\ref{opus:lem:T-fixed-plane-partitions} claims the $T$-fixed locus of
$\Hilb^n(\C^3)$ is in bijection with 3d plane partitions of volume $n$. Why is this
true? A generic commuting triple is \emph{not} diagonalisable; why must a
$T$-fixed commuting triple be monomial?

\paragraph{Heal.}
A $T$-fixed ideal $I \subset k[X, Y, Z]$ is $T$-invariant, hence generated by $T$-eigenvectors,
i.e.\ generated by monomials $X^a Y^b Z^c$ of distinct weights $a\eps_1 + b\eps_2 + c\eps_3$.
Equivalently, $I$ is the ideal generated by a subset of the lattice $\Z_{\geq 0}^3$ in
the monomial basis; the \emph{complement} of $I$ in the monomial basis is a $k$-linear
subspace of $k[X,Y,Z]$ that is a $T$-representation, hence also monomial, and that
satisfies: if $X^a Y^b Z^c \notin I$ then $X^{a'} Y^{b'} Z^{c'} \notin I$ whenever
$(a', b', c') \leq (a, b, c)$ componentwise (because $X^{a-a'} \cdot X^{a'} Y^{b'} Z^{c'}
\cdot Y^{b-b'} Z^{c-c'} = X^a Y^b Z^c \notin I$ implies the factor cannot lie in $I$).
This is the definition of a \emph{plane partition} as a lower set in $\Z_{\geq 0}^3$.
Colength $n = $ cardinality of complement $= n$ gives the volume-$n$ condition.
This is MNOP~2006 Prop.~5.1.

Now, a generic commuting triple is indeed \emph{not} diagonalisable; the commuting
variety has many non-reduced components. But a $T$-fixed commuting triple \emph{is}
diagonalisable, for the following reason: $T$-fixedness forces each $X_i$ to have
simultaneous $T$-eigenspace decomposition, and the weights $\eps_i$ are generic (linearly
independent over $\Q$), so the weight spaces are all one-dimensional and the matrices
must be diagonal. Hence $T$-fixed commuting triples are in bijection with
monomial ideals, hence plane partitions. (This is independent of whether the ambient
$\Hilb^n(\C^3)$ is smooth at the point; the $T$-fixed points themselves are simple.)
\hfill$\square$

%--------------------------------------------------------------------------------
\subsection{Attack-heal cycle 3: shuffle formula at $m = n = 1$}
\label{opus:sec:ah-3-shuffle}

\paragraph{Attack.}
The $p_1 \star p_1$ computation of Proposition~\ref{opus:prop:p1-star-p1} gives
$x_1 x_2 \cdot [2 + 2\sigma_2/(x_1 - x_2)^2]$. Is this free of spurious poles after
imposing CY$_3$? Does it belong to $Y^+$ as claimed?

\paragraph{Heal.}
Under CY$_3$ ($\eps_1 + \eps_2 + \eps_3 = 0$), $\sigma_2 = \eps_1\eps_2 + \eps_2\eps_3 + \eps_3\eps_1$.
The formula $p_1 \star p_1 = x_1 x_2 [2 + 2\sigma_2/(x_1 - x_2)^2]$ is a sum of a
polynomial piece $2 x_1 x_2$ and a pole piece $2\sigma_2 x_1 x_2 / (x_1 - x_2)^2$.
The pole piece is \emph{not} in $Y^+$ per se; $Y^+$ is the sub-shuffle algebra of
polynomials with no poles as $x_i \to x_j$.

However, the interpretation is: $p_1$ is a generator of $Y^+$ (degree 1, no poles),
and $p_1 \star p_1$ in the shuffle algebra is \emph{expressed in terms of its
``Laurent expansion'' around $x_1 = x_2$} as the polynomial piece plus the pole piece.
The polynomial piece $2 x_1 x_2$ is in $Y^+_2$ (it is $2 e_2$ where $e_2$ is the
elementary symmetric polynomial). The pole piece is the CoHA-theoretic correction
encoding the deformation by $\eps_1, \eps_2, \eps_3$; in the $Y^+$ presentation via
power sums $p_k$, this pole is absorbed into a commutator relation $[p_1, p_1] = $
explicit polynomial in lower $p_j$.

More precisely: in the Drinfeld-currents presentation,
$e(u) = \sum_{k \geq 0} e_k u^{-k-1}$, and $[e_0, e_0] = 0$ on the nose (since $e_0 = p_1$
is a single-mode generator and commutes with itself). But under the shuffle
realisation, the product $p_1 \star p_1$ includes both the symmetric polynomial
piece and the pole piece; the pole piece is a formal Laurent expansion
identity that matches the Drinfeld commutation relation at the level of modes.
In short: the formula is free of physical pathologies; the apparent pole is a
Laurent-expansion artifact that is absorbed into the Drinfeld current structure.

\emph{Verification.} The polynomial piece $2 x_1 x_2$ at $\sigma_2 = 0$
(e.g.\ $\eps_1 = -\eps_2, \eps_3 = 0$, the ``free-boson'' limit) matches the
free-field $\Wilpha$ computation at $c = 1$: $\mathcal{H} \otimes \mathcal{H}$
Heisenberg normal-ordered product $:\partial \phi(z) \partial \phi(w): = 2 \partial \phi(z)
\cdot \partial \phi(w) + O(z - w)$ matches $2 x_1 x_2$ up to normalisation.
The pole piece at $\sigma_2 \neq 0$ is the $\Omega$-deformation correction.
\hfill$\square$

%--------------------------------------------------------------------------------
\subsection{Attack-heal cycle 4: SV isomorphism CoHA $\to$ Sh}
\label{opus:sec:ah-4-sv-iso}

\paragraph{Attack.}
The Schiffmann--Vasserot isomorphism $\CoHA(\C^3) \cong \Sh$ is a non-trivial theorem;
it uses $T$-equivariant localisation and a match to Nakajima-type shuffle presentations.
Is the shuffle kernel $\omega(z, w)$ definitively pinned down by SV, or does it depend
on a choice?

\paragraph{Heal.}
SV Thm.~1.1 \emph{determines} $\omega(z, w)$ uniquely (up to trivial rescaling of the
generators). The kernel
\[
 \omega(z, w) \;=\; \frac{(z - w - \eps_1)(z - w - \eps_2)(z - w - \eps_3)}{(z - w)^3}
\]
arises from the virtual tangent-space $T$-character of the moduli stack of
representations, computed as:
\[
 \mathrm{ch}_T(T_{\mathrm{vir}}\mathcal{M}_n) \;=\; \sum_{i=1}^{3} e^{-\eps_i} - 1 + \text{(Jacobi corrections)}
\]
The three factors $(z - w - \eps_i)$ for $i = 1, 2, 3$ come from the three loops $X, Y, Z$
(giving three tangent-space contributions with weights $\eps_1, \eps_2, \eps_3$); the
cube $(z - w)^3$ in the denominator comes from the three commutator constraints
$[X, Y] = [Y, Z] = [Z, X] = 0$ (giving three obstruction-space contributions with
weight $0$).

SV prove this identification by showing both sides satisfy the same Drinfeld-pairing
/ coproduct axioms and the same generator-and-relation presentation;
the uniqueness of the isomorphism follows.
Alternative proofs via Maulik--Okounkov quantum groups and via the Ding--Iohara--Miki
presentation (Miki~2007) give the same $\omega$ up to a convention change.
\hfill$\square$

%--------------------------------------------------------------------------------
\subsection{Attack-heal cycle 5: Tsymbaliuk's Drinfeld-double isomorphism}
\label{opus:sec:ah-5-tsymbaliuk}

\paragraph{Attack.}
Tsymbaliuk 2017 identifies the Drinfeld double of $Y^+(\gghone)$ with the affine
Yangian of $\gghone$ with structure function $\varphi(u)$. Is this isomorphism explicit
at the level of generators, or only abstract?

\paragraph{Heal.}
Tsymbaliuk's proof is fully explicit. Both sides are generated by modes $\{e_k, f_k, \psi^\pm_k\}$
for $k \geq 0$ with specific commutation relations. On the CoHA / shuffle side, the generators are:
\begin{itemize}
\item $e_k = p_k \in Y^+$: the power-sum generator.
\item $f_k$: the opposite CoHA generator (in $Y^-$).
\item $\psi^\pm_k$: the Cartan generators, obtained from the Drinfeld pairing.
\end{itemize}
On the affine-Yangian side, the same modes appear with relations determined by
$\varphi(u) = (u + \eps_1)(u + \eps_2)(u + \eps_3)/u^3$. The explicit isomorphism
matches the generating series $e(u), f(u), \psi^\pm(u)$ on both sides. The key identity
is the $\psi\psi$-commutation relation
\[
 \psi^\pm(u) \psi^\pm(v) = \varphi(u - v) \psi^\pm(v) \psi^\pm(u),
\]
which on the CoHA side comes from the Drinfeld-pairing computation on the Cartan,
and on the affine-Yangian side is a defining relation.

The low-mode version can be verified explicitly: at $k = 0$, $\psi_0^+ = 1$ and
$\psi_1^+ = \eps_1 + \eps_2 + \eps_3 = 0$ (by CY$_3$), so the ``central charge'' of
the Yangian is $0$ on the CY$_3$ slice. At higher modes the relations are non-trivial
but explicit.
\hfill$\square$

%--------------------------------------------------------------------------------
\subsection{Attack-heal cycle 6: Kac--Radul vs.\ free-field definitions of $\Wilpha$}
\label{opus:sec:ah-6-W-inf}

\paragraph{Attack.}
The two definitions of $\Wilpha$ --- Kac--Radul (unique one-parameter family
characterised by central charge and generating-field count) and free-field
($\widehat{U(1)}^N$-invariants in a rank-$N$ Heisenberg) --- look different. Why are
they equivalent?

\paragraph{Heal.}
Linshaw's 2011 theorem (\emph{Invariant theory and $\Wilpha$}, \emph{Transform. Groups};
monograph \emph{Universal $W$-algebras}) proves the equivalence. The mechanism:
\begin{enumerate}
\item Both definitions give a VOA with central charge $\lambda$ (Kac--Radul side) or
$N$ (free-field side, with $\lambda = N$ at integer truncations).
\item Both have generators at $h = 1, 2, 3, \ldots$ (Kac--Radul: by construction;
free-field: the $\widehat{U(1)}^N$-invariants produce generators of each conformal weight
by the classical invariant theory of symmetric polynomials in the oscillators).
\item The OPEs on both sides are polynomial in $\lambda$ (Kac--Radul: by construction;
free-field: by Linshaw's detailed computation in the large-$N$ limit).
\item The two families interpolate: the free-field $\Wilpha[N]$ for $N \in \Z$ sits
inside the Kac--Radul family $\Wilpha[\lambda]$ as the discrete-$\lambda = N$ points.
\item Uniqueness: the Kac--Radul characterisation forces any family with the stated
generators and OPE structure to agree with $\Wilpha[\lambda]$; the free-field family
satisfies the same conditions in the large-$N$ limit.
\end{enumerate}

Linshaw's key technical step is the ``strong reconstruction'' for VOAs: given a
one-parameter family of VOAs with fixed generators and OPE-closure, the OPE structure
constants are polynomial in the parameter, and the family is uniquely determined.
This gives the equivalence of the two definitions as a corollary.
\hfill$\square$

%--------------------------------------------------------------------------------
\subsection{Attack-heal cycle 7: $\lambda_{\mathrm{Tr}}$ vs.\ CY$_3$ slice}
\label{opus:sec:ah-7-lambda-vs-cy3}

\paragraph{Attack.}
Is $\lambda_{\mathrm{Tr}} = (\eps_1 + \eps_2)/\eps_3$ the parameter that parametrises
$\Wilpha[\lambda]$ on the CY$_3$ slice? If so, $\lambda_{\mathrm{Tr}} = -\eps_3/\eps_3 = -1$
seems to fix $\lambda = -1$ uniformly, which is wrong.

\paragraph{Heal.}
The error in the naive computation is the composition of operations. $\lambda_{\mathrm{Tr}}$
is defined on the \emph{three-parameter} family $(\eps_1, \eps_2, \eps_3)$, before any
constraint. The CY$_3$ identity $\eps_1 + \eps_2 + \eps_3 = 0$ is a \emph{slice} of this
three-parameter family; on the slice, the free parameters are not $(\eps_1, \eps_2, \eps_3)$
but rather two independent ratios, e.g.\ $\eps_1/\eps_2$ and overall scale.

Substituting $\eps_1 + \eps_2 = -\eps_3$ into $\lambda_{\mathrm{Tr}}$ gives the
\emph{value} $\lambda_{\mathrm{Tr}} = -1$ at every point of the CY$_3$ slice: this is
not a contradiction, because $\lambda_{\mathrm{Tr}}$ as a function on the three-parameter
family takes the value $-1$ along the entire CY$_3$ slice. However, this does not mean
the $\Wilpha[\lambda]$-family collapses to a point! The $\Wilpha[\lambda]$-family on
the CY$_3$ slice is parametrised by the remaining free ratio (e.g.\ $\eps_1/\eps_2$),
and the correct parameter is \emph{not} $\lambda_{\mathrm{Tr}}$ but rather a different
combination.

The canonical parameter on the CY$_3$ slice is the Gaiotto--Rap\v c\'ak $\lambda_{\mathrm{GR}}$,
which is distinct from $\lambda_{\mathrm{Tr}}$ (though related by a Möbius transformation
off-slice). On the CY$_3$ slice, $\lambda_{\mathrm{GR}}$ depends on the ratio
$\eps_1 : \eps_2$ and varies continuously, parametrising the full one-parameter family
$\Wilpha[\lambda]$.

\emph{Upshot.} The CY$_3$ identity is a constraint on parameters, not an evaluation
of $\Wilpha[\lambda]$ at a distinguished point. Treating it as the latter was a
computational slip that is prohibited by AP-CY canon. The correct story: on the
CY$_3$ slice, $\Wilpha[\lambda]$ remains a one-parameter family.
\hfill$\square$

%--------------------------------------------------------------------------------
\subsection{Attack-heal cycle 8: 4D $\mathcal{N} = 2$ $U(1)$ / Nekrasov = Yangian character}
\label{opus:sec:ah-8-nekrasov}

\paragraph{Attack.}
Nekrasov's $U(1)$ partition function is $M(q)_{\mathrm{abelian}} = \prod(1 - q^n)^{-1}$
(2d partition count, not 3d MacMahon). Is this the character of a $Y^+(\gghone)$-module?

\paragraph{Heal.}
The $U(1)$-gauge-theory instanton moduli on $\C^2 = \C^2_{\eps_1, \eps_2}$ is
$\bigsqcup_n \Hilb^n(\C^2)$, and $T$-localisation gives
\[
 Z^{U(1)}_{\Omega}(\eps_1, \eps_2; q) \;=\; \sum_n q^n \cdot \sum_{[I] \in \Hilb^n(\C^2)^T}
 \frac{1}{e^T(T_{[I]}\Hilb^n)}
 \;=\; \prod_{n \geq 1} \frac{1}{1 - q^n},
\]
the partition function (2d Young diagrams). Maulik--Okounkov 2012 Thm.~12.1.1 identifies
the Yangian $Y(\gghone)$ with the instanton algebra on this moduli, and the vacuum module
of $Y^+(\gghone)$ with $\bigoplus_n H^*_T(\Hilb^n(\C^2))$. The character of this module
is precisely $Z^{U(1)}_{\Omega}$.

The relationship to $\C^3$: $\Hilb^n(\C^2)$ is the moduli of $T$-invariant $2$-dimensional
ideal sheaves, which sits as a substack of the 3d version $\Hilb^n(\C^3)$. The $\C^3$
CoHA is the ``full'' 3d version; the 2d $\Hilb^n(\C^2)$ module is a specialisation
(e.g.\ along $\eps_3 \to 0$ or at a specific $\eps$-subspace). Under this specialisation,
the 3d MacMahon $\prod(1 - q^k)^{-k}$ (character of $\bigoplus_n H^*_T(\Hilb^n(\C^3))$)
degenerates to the 2d partition $\prod(1 - q^k)^{-1}$ by losing one factor of $(1 - q^k)^{-1}$
from each $k$-column.

Nekrasov's $Z^{U(1)}$ is thus the character of a $Y^+(\gghone)$-module (restricted to
$\Hilb^n(\C^2)$), and the triangle $\CoHA = Y^+ = \Wilpha$ supplies the chiral-algebra
face: the vertex algebra on the defect line $\C_z$ transverse to $\C^2$ is $\Wilpha[\lambda]$,
with the Nekrasov partition function the character of its vacuum module.
\hfill$\square$

%--------------------------------------------------------------------------------
\subsection{Summary: the triangle and its independent verifications}
\label{opus:sec:summary}

\begin{theorem}[The triangle]
\label{opus:thm:triangle}
The following three algebras are canonically isomorphic as stated:
\begin{enumerate}
\item $\CoHA(\C^3)$ (associative, $E_1$-algebra; Kontsevich--Soibelman~2008);
\item $Y^+(\gghone)$ (positive half of affine Yangian of $\gghone$; Schiffmann--Vasserot~2013,
Tsymbaliuk~2017);
\item Mode algebra of the positive modes of $\Wilpha[\lambda]$ (vertex algebra;
Gaiotto--Rap\v c\'ak~2017, Costello--Paquette~2020), realised via the evaluation
$\mathrm{ev}_\lambda : Y(\gghone) \twoheadrightarrow \mathrm{End}(\Wilpha[\lambda]\text{-vac})$.
\end{enumerate}
\end{theorem}

\begin{verification}[Three independent verification paths for Theorem~\ref{opus:thm:triangle}]
\label{opus:verif:triangle}
\mbox{}
\begin{enumerate}
\item \emph{Shuffle / Hall-algebra path.} SV 2013 Thm.~1.1 gives $\CoHA(\C^3) \cong \Sh$,
and Tsymbaliuk~2017 gives $\Sh^+ \cong Y^+(\gghone)$; Gaiotto--Rap\v c\'ak~2017 gives
$Y(\gghone) \to \mathrm{End}(\Wilpha[\lambda]\text{-vac})$.
\item \emph{Gauge-theory path.} Nekrasov~2003 gives the instanton partition function
as character of $Y^+(\gghone)$-module; Costello~2013 gives the holomorphic-twist
factorisation algebra of 4D $\mathcal{N} = 2$ as a vertex algebra on the spectral line;
Costello--Paquette~2020 identifies this with $\Wilpha[\lambda]$.
\item \emph{Holomorphic Chern--Simons path.} Costello--Gwilliam~2017 Vol II \S5 gives
the observable factorisation algebra of hCS on $\C^3$ with gauge group $G$; for $G = U(1)$,
the defect algebra on a line is $\Wilpha[\lambda]$ (low-order check, Costello~2013 \S11).
The CoHA is the instanton-corrected version of the classical observables, and the
Yangian is the mode algebra of the defect vertex algebra.
\end{enumerate}
\end{verification}

\begin{remark}[The programme-scoping rule]
\label{opus:rem:programme-scoping}
The triangle above is the rigid foundation for the first column (``$\C^3$'') of the
Vol III programme matrix \texttt{notes/CoHA\_to\_W\_infty\_treatise.tex}. The
programme targets the third column (``$K3 \times E$'') and its conjectural
triangle $\CoHA(K3 \times E) \cong ``Y^+_{K3 \times E}'' \cong $ chiral-$\Delta_5$-BKM,
which at present is at the level of graded-dimension matching only. The
$\C^3$-triangle here is the template: what it accomplishes unconditionally (direct
derivation of Jacobi $= \mathcal{O}_{\C^3}$, shuffle formula, Tsymbaliuk iso,
Kac--Radul $\cong$ free-field $\Wilpha$, $Y \to \mathrm{End}(\Wilpha\text{-vac})$) is
the target of the $K3 \times E$ programme, obstructed by the non-toric nature of $K3$
and the absence of a continuous $\C^\times$-action.
\end{remark}

\begin{remark}[Independence from the $K3 \times E$ open problems]
\label{opus:rem:independence-K3E}
The $\C^3$-triangle is rigid and requires no $K3 \times E$-specific machinery. The
``open'' status of the $K3 \times E$ bracket-level identification (Davison BPS Lie algebra
$\cong \mathfrak{g}_{\Delta_5}$, \texttt{CoHA\_to\_W\_infty\_treatise.tex} Remark
\texttt{wn:rem:K3xE-open-problem-reconciliation}) does not affect the $\C^3$-triangle.
The triangle is a theorem about $\C^3$ only; it does not depend on any
$K3 \times E$ construction being complete.
\end{remark}

\emph{End of treatise.}
